\section{Training Details}
\label{app:details}

\paragraph{SFT stage.}
We fine-tune Qwen3-32B using LoRA with rank 64 and alpha 128 on all
linear layers. The SFT dataset consists of 5,000 samples: 637 Chinese
mathematical critique examples and 4,363 English math reasoning problems
from NuminaMath-CoT and MetaMathQA. Training runs for 3 epochs with
learning rate $5 \times 10^{-5}$, cosine decay, batch size 2, gradient
accumulation 8, and max sequence length 8,192. Total training: 120 steps,
36 minutes on 8$\times$ NVIDIA L20X. Final loss: 0.311, token accuracy:
89.7\%.

\paragraph{GRPO stage.}
We train with the TopoPRM composite reward on 637 Chinese mathematical
critique prompts. Each prompt generates $G = 8$ completions with
temperature 0.7, top-$p$ 0.95, and max completion length 2,048. Training
runs for 1 epoch with learning rate $5 \times 10^{-6}$, cosine decay,
batch size 1, and gradient accumulation 4. vLLM (colocate mode) is used
for efficient generation with $\mathit{gpu\_memory\_utilization} = 0.5$
and $\mathit{max\_model\_len} = 4{,}096$.

\paragraph{Reward weights.}
Default composite weights:
$w_o = 0.4$ (outcome),
$w_t = 0.2$ (topology),
$w_c = 0.15$ (continuity),
$w_f = 0.15$ (format),
$w_l = 0.1$ (length).

\paragraph{Hardware.}
All experiments run on 8$\times$ NVIDIA L20X GPUs (143\,GB VRAM each),
700\,GB host RAM, using the \texttt{ms-swift} 4.0.1 framework with
PyTorch 2.10.0, vLLM 0.17.1, and Transformers 4.57.6.

\section{SCAE Algorithm}
\label{app:scae}

Algorithm~\ref{alg:scae} presents the full procedure of Stratified
Clipping Advantage Estimation.

\begin{algorithm}[H]
\caption{Stratified Clipping Advantage Estimation (SCAE)}
\label{alg:scae}
\begin{algorithmic}[1]
\REQUIRE Group of completions $\{c^{(1)}, \ldots, c^{(G)}\}$ for prompt $q$
\REQUIRE Outcome rewards $\{R_\mathrm{out}^{(i)}\}_{i=1}^G$, auxiliary rewards $\{r_\mathrm{aux}^{(i)}\}_{i=1}^G$
\STATE $\bar{r}_\mathrm{acc} \leftarrow \frac{1}{G}\sum_{i=1}^G R_\mathrm{out}^{(i)}$
\STATE $\mathcal{C} \leftarrow \{i : R_\mathrm{out}^{(i)} = 1\}$ \hfill \COMMENT{correct set}
\STATE $\mathcal{W} \leftarrow \{j : R_\mathrm{out}^{(j)} = 0\}$ \hfill \COMMENT{incorrect set}
\STATE $\bar{r}_\mathrm{aux}^+ \leftarrow \mathrm{mean}(\{r_\mathrm{aux}^{(i)} : i \in \mathcal{C}\})$ if $\mathcal{C} \neq \emptyset$ else $0$
\STATE $\bar{r}_\mathrm{aux}^- \leftarrow \mathrm{mean}(\{r_\mathrm{aux}^{(j)} : j \in \mathcal{W}\})$ if $\mathcal{W} \neq \emptyset$ else $0$
\FOR{$i \in \mathcal{C}$}
    \STATE $A^{(i)} \leftarrow (1 - \bar{r}_\mathrm{acc}) + \max(0,\; r_\mathrm{aux}^{(i)} - \bar{r}_\mathrm{aux}^+)$
\ENDFOR
\FOR{$j \in \mathcal{W}$}
    \STATE $A^{(j)} \leftarrow -\bar{r}_\mathrm{acc} + \min(0,\; r_\mathrm{aux}^{(j)} - \bar{r}_\mathrm{aux}^-)$
\ENDFOR
\RETURN Advantages $\{A^{(i)}\}_{i=1}^G$
\end{algorithmic}
\end{algorithm}

\section{DAG Extraction Details}
\label{app:dag}

\paragraph{Step segmentation.}
We identify reasoning step boundaries using the following markers:
\begin{itemize}[leftmargin=*,itemsep=1pt]
    \item Numbered steps: patterns like ``Step 1:'', ``(1)'', ``第一步''
    \item Mathematical delimiters: ``$\because$'', ``$\therefore$'',
          ``$\Rightarrow$''
    \item Logical connectives: ``因为'', ``所以'', ``由此可得'', ``therefore'',
          ``hence''
    \item Line breaks with significant indentation changes
\end{itemize}

\paragraph{Dependency detection.}
For each pair of steps $(r_i, r_j)$ with $i < j$, we check:
\begin{enumerate}[leftmargin=*,itemsep=1pt]
    \item \textbf{Expression overlap}: Whether $r_j$ contains mathematical
    expressions (extracted via regex) that first appeared in $r_i$.
    \item \textbf{Variable reference}: Whether $r_j$ uses variables or
    intermediate results defined in $r_i$.
    \item \textbf{Sequential proximity}: Consecutive steps $(r_i, r_{i+1})$
    receive weak sequential edges (weight 0.3).
\end{enumerate}

\paragraph{Step classification.}
Each node is classified by keyword matching into one of seven types:
\textsc{definition} (``设'', ``let'', ``定义''),
\textsc{derivation} (``由'', ``根据'', ``since''),
\textsc{computation} (``计算'', ``代入'', ``=$\cdots$=''),
\textsc{conclusion} (``因此'', ``所以'', ``therefore''),
\textsc{auxiliary} (``注意'', ``note''),
\textsc{substitution} (``代入'', ``substituting''),
\textsc{case\_analysis} (``分类讨论'', ``case'').

\section{Case Study}
\label{app:case_study}

We present a representative example comparing the reasoning structure
of completions from the SFT baseline, outcome-only GRPO, and TopoPRM
(full) on a middle-school algebra problem. The case study will be
populated with actual model outputs after training completes.

\paragraph{Problem.}
A middle-school algebra problem requiring multi-step equation solving
with verification of intermediate steps.

\paragraph{Analysis.}
The TopoPRM-trained model is expected to produce reasoning traces with:
(i) clearly delineated steps forming a valid DAG,
(ii) each conclusion explicitly referencing its supporting premises,
(iii) more concise overall length compared to outcome-only GRPO.
In contrast, the outcome-only model may produce correct answers but with
less structured, potentially redundant reasoning chains.

\section{Distillation Details}
\label{app:distill}

For reasoning compression, we use the TopoPRM-trained Qwen3-32B as the
teacher model to generate high-quality reasoning traces. We filter
teacher outputs to retain only those with composite reward
$R_\mathrm{total} > 0.7$. The student models (Qwen3-14B, Qwen3-7B) are
then trained using reverse-KL distillation with LoRA (rank 32, alpha 64)
for 3 epochs.
